\section{Programmation par continuation}

L'outil principal permettant la conversion de fonctions récursives quelconques
en fonctions récursives terminales est la
\textit{programmation par continuation} ou CPS (de son nom anglais
\textit{continuation passing style}). Toute fonction écrite sous ce paradigme de
programmation prend un argument supplémentaire, une \textit{continuation}, qui
est une fonction à un argument du type du résultat de la fonction. Cette
continuation est alors appelée avec le résultat une fois celui-ci calculé. Par
exemple, la fonction \texttt{sum} précédente peut être ré-écrite sous CPS comme
suit :

\small
\begin{minted}{ocaml}
    let rec sum_cps (l : int list) (continuation : int -> 'a') : 'a =
        match l with
        | [] -> continuation 0
        | x :: q ->
            let faire_addition (v : int) : 'a =
                continuation (x + v)
            in
            sum_cps q faire_addition

    let sum (l : int list) : int = sum_cps l (fun res -> res)
\end{minted}
\normalsize

Ici, la continuation est en effet appelée avec \texttt{0} lorsque
\texttt{sum\_cps} est appelée sur la liste vide, puisque la somme des éléments
de la liste vide est nulle. Lorsque \texttt{l} est non vide, une nouvelle
continuation \texttt{faire\_addition} est créée : celle-ci prend alors le
résultat de l'appel récursif (dans la deuxième branche du \texttt{match}), y
ajoute la valeur de \texttt{x}, puis appelle la continuation originale avec la
nouvelle valeur. Cela correspond bien à la somme de tous les éléments de
\texttt{l}. Pour obtenir le résultat, il suffit d'appeler la fonction sous CPS
avec l'identité comme continuation.

On peut de plus remarquer que cette implémentation de \texttt{sum\_cps}
est récursive terminale (et la nouvelle continuation introduite aussi). Notre
algorithme aura alors pour but de convertir une fonction récursive quelconque en
une fonction sous CPS, dont nous pourrons montrer qu'elle est récursive
terminale.

Notre algorithme est alors constitué de trois étapes principales :
\begin{enumerate}
    \item identifier les appels récursifs~;
    \item mettre les opérations survenant après un appel récursif dans une
          continuation prenant en argument le résultat de cet appel récursif~;
    \item modifier le programme pour le mettre sous CPS avec les continuations
          créées.
\end{enumerate}
